\chapter{Физические углы и классификация неоднозначности связности}

Связывающий родитель Мориты определил носитель переходов и двустороннюю
кривизну. Остался решающий вопрос: выделяют ли известные семейное и
стандарт-модельное действия единственную юкавскую компоненту связности?

\section{Антикруговая сверка}

В Томе IV уже были установлены два запрета. Один и тот же набор семейных
данных допускал несколько неэквивалентных ковариантных юкавских карт, а
обычная внутренняя флуктуация либо не содержала требуемое направление,
либо вместе с ним восстанавливала всё $M_3(\mathbb C)$. Новый связывающий
контейнер не считается решением, пока он не устраняет именно эту
неединственность.

Литературная граница согласуется с такой постановкой: конечный оператор
Дирака некоммутативной Стандартной модели кодирует юкавские связи и их
параметры, тогда как диаграмма Краевского ограничивает допустимый рисунок
блоков. См. \path{arXiv:hep-th/0608226}, \path{arXiv:0901.0577} и
\path{arXiv:0809.5137}. Поэтому требуется отдельная теорема выбора, а не
повторное перечисление разрешённых рёбер.

\section{Разность двух связностей}

Пусть две двусторонние связности на одном бимодуле имеют одинаковые левое
и правое правила Лейбница. Их разность является бимодульно-линейным
эндоморфизмом. Для полного бимодуля
\begin{equation}
 {}_{M_{20}}E_{M_{15}}=M_{20\times15}(\mathbb C)
\end{equation}
двойная теорема о коммутанте даёт
\begin{equation}
 \operatorname{End}_{M_{20}-M_{15}}(E)=\mathbb C I_E.
 \label{eq:v5-full-factor-connection-ambiguity}
\end{equation}
То есть полные матричные углы оставляют только общий скаляр. Но правый
$M_{15}$ был контрольным чтением, а не физической координатной алгеброй.

\section{Ограничение до наблюдаемых блоков}

Точное пятнадцатимерное чтение разбито на блоки
\begin{equation}
 H_{15}=Q_L^{(6)}\oplus L_L^{(2)}\oplus u_R^{(3)}
 \oplus d_R^{(3)}\oplus e_R^{(1)}.
\end{equation}
Максимальная сохраняющая это разложение алгебра равна
\begin{equation}
 \mathcal B_{\mathrm{obs}}
 =M_6(\mathbb C)\oplus M_2(\mathbb C)
 \oplus M_3(\mathbb C)\oplus M_3(\mathbb C)\oplus\mathbb C.
\end{equation}
Её коммутант имеет вид
\begin{equation}
 \mathcal B_{\mathrm{obs}}'
 =\mathbb CP_Q\oplus\mathbb CP_L\oplus\mathbb CP_u
 \oplus\mathbb CP_d\oplus\mathbb CP_e,
 \qquad \dim_{\mathbb C}=5.
 \label{eq:v5-observed-block-commutant}
\end{equation}
После удаления общего скаляра остаётся четырёхмерное комплексное
пространство.

Физическое представление
$\mathbb C\oplus\mathbb H\oplus M_3(\mathbb C)$ содержится в
$\mathcal B_{\mathrm{obs}}$. Поэтому его коммутант не меньше
\eqref{eq:v5-observed-block-commutant}. Число пять является нижней границей,
а не следствием слишком широкого контрольного представления.

Машинная классификация всех $225$ матричных направлений подтверждает
\begin{equation}
 \dim M_{15}'=1,
 \qquad
 \dim\mathcal B_{\mathrm{obs}}'=5
\end{equation}
с максимальным дефектом коммутатора $2.3\cdot10^{-16}$.

\section{Почему ранг-один семейный угол не спасает единственность}

Пусть $P_\rho$ --- уже выведенный семейный проектор ранга один. Ограничение
носителя до $P_\rho E$ действует только слева. Каждое правое умножение на
элемент $\mathcal B_{mathrm{obs}}'$ по-прежнему сохраняет этот угол:
\begin{equation}
 P_\rho(\xi C)=(P_\rho\xi)C,
 \qquad C\in\mathcal B_{\mathrm{obs}}'.
\end{equation}
Поэтому после выбора семейного состояния остаются те же пять наблюдаемых
направлений связности. Машинные дефекты лево-правой ассоциативности и её
углового ограничения равны нулю.

Этот результат уточняет роль ранга один. Он выбирает семейный конец
перехода, но не выбирает, в какой наблюдаемый блок должен входить переход и
с каким относительным оператором.

\section{Что фиксирует единственный след}

Наблюдаемый след задаёт диагональную метрику на пяти проекторах:
\begin{equation}
 G_{ss'}=\tau_{15}(P_sP_{s'})
 =\delta_{ss'}\left(\frac25,\frac2{15},\frac15,
 \frac15,\frac1{15}\right)_s.
 \label{eq:v5-observed-connection-metric}
\end{equation}
Тем самым нормы и кратности фиксированы. Но метрика не выбирает ненулевой
вектор в четырёхмерном центрированном пространстве. Чистый положительный
квадратичный функционал имеет единственный минимум в нулевой связности, то
есть уничтожает, а не выводит юкавское взаимодействие.

\begin{theorem}[Неединственность физической связности]
Связывающий родитель $M_{35}$ однозначно фиксирует носитель переходов,
левое и правое действия, знак относительной кривизны и метрику следа. После
ограничения наблюдаемого угла до физического блочного чтения пространство
разностей совместимых связностей имеет комплексную размерность не менее
пяти, или не менее четырёх после удаления общего скаляра. Ранг-один
семейный проектор эту неоднозначность не уменьшает.
\end{theorem}

\begin{nogocriterion}[Запрет объявления юкавского вывода]
Нельзя считать юкавский оператор выведенным только из связывающей алгебры,
условных ожиданий и ранг-один проектора. Эти данные определяют место и
метрику взаимодействия, но не выбирают единственную ненулевую компоненту
связности.
\end{nogocriterion}

\section{Оставшийся динамический тест}

Повторный поиск алгебры или проекции запрещён прежними гейтами. Допустим
только один следующий вопрос: выбирает ли уже существующая родительская
кривизна ненулевую единственную орбиту в четырёхмерном центрированном
пространстве наблюдаемых компонентов связности?

Нужно выписать самый общий функционал, разрешённый текущими симметриями и
единственным следом, без подстановки масс. Если его минимум нулевой,
вырожденный или содержит свободные коэффициенты, происхождение юкавского
оператора следует заморозить как недоопределённое.

\begin{protocol}
\begin{enumerate}
 \item полные матричные углы: \textbf{неоднозначность один скаляр};
 \item физическое наблюдаемое ограничение:
 \textbf{не менее пяти комплексных направлений};
 \item центрированная неоднозначность:
 \textbf{не менее четырёх комплексных направлений};
 \item ранг-один семейный угол как селектор наблюдаемого блока:
 \textbf{не проходит};
 \item единственный след как метрика: \textbf{проходит};
 \item единственный след как селектор ненулевой связности:
 \textbf{не проходит};
 \item связывающий родитель Мориты: \textbf{сохранён};
 \item единственный юкавский оператор: \textbf{не выведен};
 \item физическое замыкание: \textbf{не достигнуто}.
\end{enumerate}
\end{protocol}

Машинный сертификат сохранён в
\path{s2t_v5_physical_corner_connection_classification_gate.py} и
\path{s2t_v5_physical_corner_connection_classification_gate_results.json}.